%%%%%%%%%%%%%%%%%%%%%part-03.tex%%%%%%%%%%%%%%%%%%%%%%%%%%%%%%%%%%
% 
% sample part title
%
% Use this file as a template for your own input.
%
%%%%%%%%%%%%%%%%%%%%%%%% Springer %%%%%%%%%%%%%%%%%%%%%%%%%%

\begin{partbacktext}
\part{Fenómenos emergentes y lógica de colisiones}

\noindent Una vez establecida la base axiomática en la Parte I y formalizada la infraestructura geométrica en la Parte II, esta tercera sección del documento se centra en la ejecución y evaluación del modelo computacional. El objetivo primordial es demostrar cómo la dinámica de los sistemas complejos, regida por reglas de transición locales, permite la emergencia de estructuras móviles capaces de procesar información de manera determinista.

En el Capítulo \ref{cap:dinamica_particulas}, se aborda el estudio cinemático de las partículas móviles o \textit{gliders}. Se define el marco algorítmico necesario para el rastreo de estas estructuras en la malla transformada, analizando su comportamiento bajo las reglas \textit{Spiral} y \textit{Beehive}. A través de la implementación de un motor de simulación eficiente, se verifica la síntesis de operadores booleanos universales mediante la lógica de colisiones, cumpliendo con los criterios sintéticos establecidos originalmente en la Sección \ref{sec:objetivos}.

Posteriormente, se procede a la validación cuantitativa y cualitativa del sistema. En el Capítulo 6, se definen métricas basadas en la teoría de la información y la entropía para evaluar el estado de la red durante su evolución. Se realiza un análisis asintótico de las estructuras emergentes para determinar sus tiempos de estabilización y ciclos límite, garantizando la robustez del isomorfismo teórico desarrollado en el Capítulo \ref{chap:isomorfismos} frente a la experimentación práctica.

Finalmente, el Capítulo 7 consolida los hallazgos de la investigación. Se presenta una síntesis analítica que contrasta las hipótesis iniciales con los datos obtenidos del simulador en C/C++. En este cierre, se discuten los límites del mapeo topológico, el costo computacional de la transformación de vecindades y se proponen nuevas líneas de investigación enfocadas en la generalización de estos modelos hacia espacios celulares de mayor complejidad.

\end{partbacktext}